\documentclass[12pt,a4paper]{article}
\usepackage[utf8]{inputenc}
\usepackage[german]{babel}
\usepackage[T1]{fontenc}
\usepackage{newtxtext,newtxmath}
\usepackage{microtype}
\usepackage{enumitem}
\usepackage{longtable,array}
\usepackage{listings,xcolor}
\usepackage[hidelinks]{hyperref}
\usepackage{xurl}
\hypersetup{pdftitle={Programmierhandbuch tails-pdp},pdfauthor={Christian Koch}}

\usepackage[left=4cm,right=2cm,top=3cm,bottom=3cm,footskip=1.75cm]{geometry}

% FTK-Layout: einzeiliger Text mit einer Leerzeile zwischen Absätzen.
\usepackage[skip=\baselineskip,indent=0pt]{parskip}
\usepackage{titlesec}
\titleformat{\section}{\normalfont\raggedright\fontsize{16}{19.2}\selectfont\bfseries}{\thesection}{1em}{}
\titleformat{\subsection}{\normalfont\raggedright\fontsize{14}{16.8}\selectfont\bfseries}{\thesubsection}{1em}{}
\titleformat{\subsubsection}{\normalfont\raggedright\fontsize{12}{14.4}\selectfont\bfseries}{\thesubsubsection}{1em}{}
\titleformat{\paragraph}[block]{\normalfont\raggedright\fontsize{12}{14.4}\selectfont\itshape}{\theparagraph}{1em}{}
\titlespacing*{\section}{0pt}{0pt}{\baselineskip}
\titlespacing*{\subsection}{0pt}{2\baselineskip}{\baselineskip}
\titlespacing*{\subsubsection}{0pt}{2\baselineskip}{\baselineskip}
\titlespacing*{\paragraph}{0pt}{2\baselineskip}{\baselineskip}

% Tabellen und Beschriftungen einheitlich in 10 pt.
\usepackage{etoolbox}
\AtBeginEnvironment{tabular}{\fontsize{10}{12}\selectfont}
\usepackage{caption}
\DeclareCaptionFont{ftk}{\fontsize{10}{12}\selectfont}
\captionsetup{font=ftk,format=hang,justification=raggedright,singlelinecheck=false}
\captionsetup[table]{position=top}
\captionsetup[figure]{position=bottom}

% Seitenzahl rechts unten; Grundschrift der Fußzeile 10 pt, Seitenzahl 12 pt.
\usepackage{fancyhdr}
\fancypagestyle{ftk}{
    \fancyhf{}
    \fancyfoot[R]{\normalfont\fontsize{10}{12}\selectfont{\fontsize{12}{14.4}\selectfont\thepage}}
    \renewcommand{\headrulewidth}{0pt}
    \renewcommand{\footrulewidth}{0.4pt}
}
\pagestyle{ftk}
\setlength{\emergencystretch}{3em}

% Codebeispiele in 10 pt, passend zur Tabellenschrift der Thesis.
\AtBeginEnvironment{longtable}{\fontsize{10}{12}\selectfont}
\lstset{basicstyle=\ttfamily\fontsize{10}{12}\selectfont,
  breaklines=true,columns=fullflexible,keepspaces=true,
  showstringspaces=false,frame=single,framerule=0.4pt,
  backgroundcolor=\color{white},aboveskip=\baselineskip,belowskip=\baselineskip}

\title{Programmierhandbuch\\tails-pdp\\\large Architektur, Schnittstellen und Erweiterung}
\author{Christian Koch}
\date{Stand: 25. Oktober 2026}
\begin{document}
\hypersetup{pageanchor=false}
\pagenumbering{gobble}
\maketitle
\thispagestyle{empty}

\clearpage
\pagenumbering{roman}
\hypersetup{pageanchor=true}
\tableofcontents

\clearpage
\pagenumbering{arabic}

\section{Zweck und Quellstand}
Dieses Handbuch beschreibt die Programmierung und Wartung des ASBAC-Prototyps
\texttt{tails-pdp}. Grundlage ist der bereinigte Workspace mit Skripten unter
\path{scripts/} und ohne die stillgelegten Socket-Bind-Datentypen.
Die Installation auf dem Linux-Zielsystem ist im separaten Installationshandbuch
beschrieben. Die Dokumente sind unabhängig voneinander kompilierbar.

Die Runtime unterstützt \texttt{file\_open}: Neue Dateiöffnungen werden über BPF-LSM
entschieden. Ein Userspace-PEP bewertet bereits offene Dateideskriptoren bei
relevanten Änderungen nach. Policies und Attribute werden aus Dateien geladen;
das Admintool ist eine lesende Diagnosekomponente. Der Parser unterstützt einen
beschränkten SAPL-inspirierten Ausschnitt, keine vollständige SAPL-Engine.

\section{Workspace und Verantwortlichkeiten}
Alle Pfade beziehen sich, sofern nicht anders angegeben, auf die Repository-Wurzel.
\begin{longtable}{>{\raggedright\arraybackslash}p{0.41\textwidth} >{\raggedright\arraybackslash}p{0.53\textwidth}}
\textbf{Crate bzw. Verzeichnis} & \textbf{Aufgabe} \\
\hline
\endhead
\path{tails-pdp/} & Userspace-Einstieg, eBPF-Laden, Initialisierung, Attach und Lebenszyklus der asynchronen Aufgaben. \\
\path{tails-pdp-ebpf/} & LSM-Hook, Tail-Call-Kette, Maps und Auswertung neuer Dateiöffnungen; \texttt{no\_std}. \\
\path{tails-pdp-common/} & Gemeinsame Datentypen, Map-Layouts, Konstanten und Policy-/Attributlogik; \texttt{no\_std}, optionale Userspace-Funktionen über Feature \texttt{user}. \\
\path{tails-pdp-policy-loader/} & Dateiparser, Übersetzung und generationenbasierte Aktivierung der Policies. \\
\path{tails-pdp-attribute-loader/} & Zeitmap, strukturierte Attribute und Aktivierung neuer Attributgenerationen. \\
\path{tails-pdp-userspace-common/} & Gepinnte Maps, Aktivierungstrigger, rekursiver Dateiwatcher und Messmarkierungen. \\
\path{tails-pdp-userspace-pep/} & Nachbewertung offener FDs und versuchter Entzug verletzender Deskriptoren. \\
\path{tails-pdp-admintool/} & Kommandozeile und lesende Ausgabe von Policies, Attributen und aktivem Zustand. \\
\path{scripts/tests/} & Komponenten-, E2E-, Performance-, Last-, Stabilitäts- und Demo-Regressionsprüfungen einschließlich Pythonhelfern. \\
\path{scripts/demo/} & Geführte tmux-Demo, Offline-Buildhelfer, Werkzeugprüfung und isolierter Demoablauf. \\
\path{examples/} & Inaktive Policy- und Attributbeispiele zur gezielten Veröffentlichung. \\
\end{longtable}

\section{Kontrollfluss und Entscheidungslogik}
\subsection{Initialisierung und Lebenszyklus}
\path{tails-pdp/src/main.rs} richtet das Pin-Verzeichnis ein und prüft vorhandene
Map-Layouts. Es lädt das eingebettete eBPF-Objekt und die BTF-Daten, lädt die
LSM-Programme und belegt die Tail-Call-Tabelle. Danach werden Policies, Zeit und
Attribute initial geschrieben. Erst anschließend wird der Einstiegshook angehängt.

Nach dem Attach laufen Zeitpflege, Attributwatcher, Policysynchronisierung und
Userspace-PEP über \texttt{tokio::select!}. Strg+C beendet die Runtime; die Rückkehr
einer dieser Aufgaben beendet ebenfalls den gemeinsamen Ablauf und propagiert
gegebenenfalls den Fehler. Es gibt keinen separat betriebenen Policy-Daemon und
keine HTTP-Schnittstelle. Gepinnte Maps können das Prozessende überdauern.

\subsection{Neue Dateiöffnung}
Der Hook \texttt{file\_open} erhält Requestdaten aus dem Kernel. Die Ressource wird
über Gerät und Inode identifiziert, der Prozess über UID und Kernel-Commandnamen.
Ein bereits negativer Rückgabewert einer vorausgehenden BPF-LSM-Entscheidung wird
erhalten. Danach folgt die Kette:
\begin{lstlisting}
file_open
  -> evaluate_file_open_static_policies  (Tail-Index 0)
  -> evaluate_file_open_stream_policies  (Tail-Index 1)
  -> combine_file_open                  (Tail-Index 2)
\end{lstlisting}
Die Jump-Table-Indizes und Programmnamen sind in \path{tails-pdp/src/lib.rs} und
\path{tails-pdp-ebpf/src/maps.rs} aufeinander abgestimmt.
Ein fehlgeschlagener Tail-Call oder fehlender Entscheidungszustand führt zur
Ablehnung. Die Kombination verwendet Deny-Overrides: Eine passende Deny-Policy
führt zu \texttt{-EPERM}; ohne passende Deny-Policy ist die Entscheidung erlaubend.
Eine Permit-Policy überstimmt kein Deny. Ein leerer Policystand erlaubt daher
Dateiöffnungen, sofern die übrigen Sicherheitsmechanismen sie erlauben.

\subsection{Nachbewertung offener Deskriptoren}
\path{tails-pdp-userspace-pep/src/pep.rs} empfängt Aktivierungstrigger und plant
zusätzlich Nachbewertungen an relevanten Zeitbedingungsgrenzen. Es scannt
\path{/proc} nicht periodisch unabhängig von diesen Ursachen.
Ein Scan verwendet einen gemeinsamen Kontext aus Policygeneration,
Attributgeneration, Attributbank und Zeit. Prozess- und FD-Daten werden auf
Dateiidentitäten zurückgeführt und gegen die Policylogik geprüft.

Der Entzug wird über \path{tails-pdp-userspace-pep/src/fd_revoker.rs} versucht.
Unter x86\_64-Linux wird der Zielprozess mit \texttt{ptrace} angehalten; ein
\texttt{close}-Systemaufruf wird injiziert und Code-/Registerzustand werden
wiederhergestellt. Auf anderen Plattformen liefert die Entzugsimplementierung
einen ausdrücklichen Fehler. PID 0, PID 1 und der Runtime-Prozess sind geschützt.
Der Ablauf ist nicht atomar mit allen konkurrierenden Prozessoperationen.

\section{Maps und binäre Schnittstellen}
Die Maps werden in \path{tails-pdp-ebpf/src/maps.rs} definiert.
Gepinnte Maps liegen unter \path{/sys/fs/bpf/tails-pdp/}.
\begin{longtable}{>{\raggedright\arraybackslash}p{0.46\textwidth} >{\raggedright\arraybackslash}p{0.48\textwidth}}
\textbf{Map} & \textbf{Inhalt und Kapazität} \\
\hline
\endhead
\texttt{FILE\_OPEN\_STATIC\_POLICIES} & Array mit 32 \texttt{FileOpenStaticPolicy}-Einträgen: zwei Bänke mit je 16 Policies. \\
\texttt{FILE\_OPEN\_STREAM\_POLICIES} & Array mit 32 \texttt{FileOpenStreamPolicy}-Einträgen: zwei Bänke mit je 16 Policies. \\
\texttt{POLICY\_GENERATION} & Array mit einer aktiven Generationsnummer vom Typ \texttt{u32}. \\
\texttt{ATTRIBUTES} & HashMap mit höchstens 1024 Einträgen; Schlüssel \texttt{AttributeKey}, Wert \texttt{AttributeValue}. \\
\texttt{ATTRIBUTE\_GENERATION} & Array mit einer aktiven Attributgeneration vom Typ \texttt{u32}. \\
\texttt{CURRENT\_TIME} & Array mit einer Unix-Zeit in Sekunden vom Typ \texttt{u64}. \\
\texttt{FILE\_OPEN\_JUMP\_TABLE} & Nicht gepinnte ProgramArray mit drei Tail-Call-Zielen. \\
\texttt{DECISIONS} & Nicht gepinnte PerCpuArray mit Deny-, Permit- und Policygenerationszustand. \\
\texttt{DEBUG\_LOGGING} & Nicht gepinnte Array mit einem Debug-Schalter. \\
\end{longtable}

Die gemeinsamen \texttt{repr(C)}-Strukturen und expliziten Paddingfelder bilden
die Schnittstelle zwischen Kernel und Userspace. \texttt{repr(u8)}-Enums haben
explizite Werte; \texttt{PolicyAction::FileOpen} bleibt \texttt{1}.
Die File-Open-Policystrukturen besitzen im geprüften Layout 104 bzw. 320 Byte
bei Ausrichtung 8. Direkte Map-Zugriffe erfordern exakt passende Typen und Größen.
Padding, Feldreihenfolge und Enumwerte dürfen nicht beiläufig verändert werden.
\texttt{aya::Pod}-Implementierungen sind Teil dieser Verantwortung.

Die Policybank beginnt bei
\texttt{(generation \% 2) * 16}; die Attributbank wird durch
\texttt{generation \% 2} gewählt. Loader schreiben zunächst den inaktiven Stand
und schalten danach die Generation um. Policy- und Attributgeneration sind
getrennte Zustände; mehrere veröffentlichte Dateien und die beiden Loader sind
keine gemeinsame Transaktion. Die Attributkapazitätsprüfung berücksichtigt
die tatsächlich gepinnte Map und den noch benötigten aktiven Bestand. Eine feste
Garantie von 512 Attributen je Generation lässt sich daraus nicht ableiten.

\section{Dateibasierte Konfigurationsschnittstelle}
\subsection{Policyformat}
Unter \path{policies/} einschließlich Unterverzeichnissen wird jede Datei mit
Endung \texttt{.policy} als genau ein Policydokument gelesen. Der Pfad ergibt
sich aus dem Arbeitsverzeichnis der Runtime, nicht dem Ort des Programmbinaries.
\begin{lstlisting}
policy "deny demo for engineers"
deny
    action == "file_open";
    command == "cat";
    resource.path == "/tmp/tails-pdp-programmer-demo.txt";
    subject.position == "engineer";
    system.defcon <= 3;
\end{lstlisting}
Die Entitlement-Zeile ist \texttt{permit} oder \texttt{deny}. Statements enden mit
Semikolon und werden als Konjunktion verknüpft. Unterstützt sind:
\begin{itemize}[nosep]
  \item \texttt{action == "file\_open"} (auch Schreibweise \texttt{file-open});
  \item optionale statische Filter \texttt{subject.uid}, \texttt{command} und \texttt{resource.path};
  \item Zeitbedingungen wie \texttt{environment.time \% N < VALUE} oder
    Vergleiche mit \texttt{environment.utc.hour}, \texttt{minute} und \texttt{second};
  \item bis zu vier dynamische Bedingungen auf \texttt{system.*}, \texttt{subject.*}
    und \texttt{resource.*} sowie höchstens eine zusätzliche Zeitbedingung;
  \item numerische Vergleiche \texttt{<}, \texttt{<=}, \texttt{==}, \texttt{>=},
    \texttt{>}; Stringattribute werden nur mit Gleichheit verglichen.
\end{itemize}
Nicht unterstützt sind unter anderem Socket-Bind-Aktionen, Imports, Schemata,
Variablen, Obligations, Advice, Transformationen und \texttt{or}.
Disjunktionen können für passende Anwendungsfälle als mehrere Policies
modelliert werden; dabei gilt weiterhin Deny-Overrides.

Die maximale Commandlänge beträgt 15 Byte; Linux reserviert im 16-Byte-Commandfeld
ein Nullbyte. Ressourcenpfade sind auf 64 Byte begrenzt. Der Loader löst vorhandene
Ressourcenpfade in Gerät/Inode auf. Eine später am selben Pfad ersetzte Datei besitzt
möglicherweise eine andere Identität; die Regel ist keine laufende Pfadnamensüberwachung.

\subsection{Attributformat}
Die Dateien verwenden die Endung \texttt{.attributes} und Einträge
\texttt{name = wert}. Beispiele:
\begin{lstlisting}
# attributes/system.attributes
defcon = 3

# attributes/subjects/1000.attributes
position = "engineer"
access-level = 2
on_call = true

# attributes/resources/tmp/tails-pdp-programmer-demo.txt.attributes
classification = "internal"
clearanceLevel = 3
\end{lstlisting}
Die drei Namensräume sind System, Subject und Resource. Subjectdateien werden
über die UID benannt. Resource-Dateipfade bilden den absoluten Ressourcenpfad
ohne führenden Schrägstrich plus Attributsuffix ab; beim Laden wird die Identität
der beschriebenen Datei aufgelöst. Globale Systemattribute haben keine
requestabhängige Objektidentität. Namen und Stringwerte werden durch gemeinsame
Hashfunktionen in die Map-Repräsentation überführt. Zahlen sind vorzeichenlose
Ganzzahlen, boolesche Werte sind \texttt{true}/\texttt{false}; DEFCON muss zwischen
1 und 5 liegen. Regeln mit \texttt{subject.position} benötigen die Attributdatei
für die tatsächlich anfragende UID, nicht zwingend für UID 1000.

\subsection{Veröffentlichung und Fehlerverhalten}
\path{scripts/publish-config.sh} erhält genau zwei Argumente: eine vollständig
vorbereitete Quelldatei und das Ziel mit Endung \texttt{.policy} oder
\texttt{.attributes}. Es schreibt eine temporäre Datei im Zielverzeichnis und
ersetzt das Ziel durch Umbenennung auf demselben Dateisystem. Bei Ersetzung werden
die Zugriffsbits des vorhandenen Ziels erhalten; Eigentümer ist der veröffentlichende
Prozess. Symlink- und Verzeichnisziele werden abgelehnt.
\begin{lstlisting}
printf 'demo\n' > /tmp/tails-pdp-programmer-demo.txt
bash scripts/publish-config.sh /tmp/prepared.policy policies/demo.policy
bash scripts/publish-config.sh /tmp/prepared.attributes attributes/system.attributes
rm policies/demo.policy
\end{lstlisting}
Die vorbereiteten Quelldateien müssen vorher mit vollständigen gültigen Inhalten
erstellt werden. Der rekursive Inotify-Watcher reagiert auf Veröffentlichung und
Entfernung; temporäre Schreibvorgänge werden ignoriert. Es gibt keinen festen
Debounce-Timer. Die Veröffentlichung validiert die Policysemantik nicht und
bestätigt nicht die spätere Aktivierung. Laufende ungültige Updates erhalten den
letzten gültigen Stand; initial ungültige Daten können den Start verhindern.
Bestätigung erfolgt über Runtime-Log und Admintool.

\section{Rust-Schnittstellen und Fehlerkontrakte}
Die folgenden Schnittstellen werden vom Runtime-Einstieg verwendet. Sie ersetzen
keine vollständige Rustdoc-Auflistung, beschreiben aber die zentralen Integrationspunkte.
\subsection{Policyloader}
\begin{lstlisting}
PolicyDirectorySync::new(
    enforcement_triggers: mpsc::Sender<EnforcementTrigger>
) -> anyhow::Result<Self>

PolicyDirectorySync::directory(&self) -> &Path
PolicyDirectorySync::sync_initial(&mut self) -> anyhow::Result<()>
PolicyDirectorySync::run(self) -> impl Future<Output = anyhow::Result<()>>
verify_pinned_map_layouts() -> anyhow::Result<()>
\end{lstlisting}
\texttt{new} erstellt das Standardverzeichnis und öffnet die benötigten gepinnten
Maps. \texttt{sync\_initial} validiert und aktiviert den vollständigen initialen
Stand; Fehler werden zurückgegeben. \texttt{run} ist im Quellcode eine
\texttt{async fn}, konsumiert das Objekt und verarbeitet Verzeichnisänderungen.
Ein Aktivierungstrigger wird erst nach erfolgreichem Commit ausgelöst.
Die Layoutprüfung erfolgt vor dem Laden des eBPF-Objekts.

\subsection{Attribut- und Zeitloader}
\begin{lstlisting}
open_attribute_maps() -> anyhow::Result<AttributeMaps>
write_current_attributes(
    attribute_maps: &mut AttributeMaps,
    enforcement_triggers: &mpsc::Sender<EnforcementTrigger>
) -> anyhow::Result<()>
async fn run_attribute_updater(
    attribute_maps: &mut AttributeMaps,
    enforcement_triggers: mpsc::Sender<EnforcementTrigger>
) -> anyhow::Result<()>

open_current_time_map() -> anyhow::Result<Array<MapData, u64>>
write_current_time(map: &mut Array<MapData, u64>) -> anyhow::Result<()>
async fn run_current_time_updater(
    map: &mut Array<MapData, u64>
) -> anyhow::Result<()>
\end{lstlisting}
\texttt{AttributeMaps} kapselt Attribute und Generationsmap. Der initiale
Attributschreibvorgang lädt die Dateien und aktiviert eine Generation.
Der Updater validiert neue Inhalte und Kapazitäten vor der Aktivierung.
Die Zeitpflege schreibt Unix-Sekunden in \texttt{CURRENT\_TIME[0]} und aktualisiert
sie im Ein-Sekunden-Takt. UTC-Komponenten werden durch \texttt{PolicyTime} abgeleitet.

\subsection{Trigger und Userspace-PEP}
\begin{lstlisting}
enum EnforcementTrigger {
    PolicyGenerationActivated { generation: u32 },
    AttributeGenerationActivated { generation: u32 },
    TimeConditionChanged { unix_seconds: u64 },
}

notify_enforcement(
    sender: &mpsc::Sender<EnforcementTrigger>,
    trigger: EnforcementTrigger
) -> anyhow::Result<bool>

async fn run_userspace_pep(
    enforcement_triggers: mpsc::Receiver<EnforcementTrigger>
) -> anyhow::Result<()>
\end{lstlisting}
Der Kanal hat Kapazität 1. \texttt{notify\_enforcement} verwendet
\texttt{try\_send}: \texttt{true} bedeutet eingereiht; ein voller Kanal liefert
\texttt{false}, weil eine Nachbewertung bereits aussteht. Ein geschlossener Kanal
ist ein Fehler. Trigger sind keine verlustfreie Historie aller Generationen.
Zeittrigger werden vom PEP beim Erreichen einer berechneten Zeitgrenze intern
gebildet; sie enthalten deren Unix-Sekunden. Der PEP bewertet den aktuellen
aktiven Zustand. Der FD-Entzugshelfer ist ein
internes Modul und wird nicht als öffentliche Crate-Schnittstelle exportiert.

\subsection{Gemeinsame Policylogik}
\path{tails-pdp-common/src/lib.rs} enthält unter anderem
\texttt{FileOpenRequest}, \texttt{FileOpenStaticPolicy},
\texttt{FileOpenStreamPolicy}, \texttt{AttributeCondition}, \texttt{AttributeKey},
\texttt{AttributeValue} und \texttt{PolicyTime}.
\path{evaluate_file_open_static_policy} liefert ein optionales Entitlement;
\texttt{None} bezeichnet eine nicht passende Policy. Die Stream-Auswertung trennt
statische Requestfilter, Zeitbedingung und dynamische Map-Abfragen.
Die noch verwendeten Funktionen mit \texttt{legacy} im Namen verarbeiten aktive
Zeitbedingungen und dürfen nicht allein wegen ihrer Benennung entfernt werden.

\section{Admintool und Diagnoseaufrufe}
Das Admintool öffnet Maps lesend. Es bietet \texttt{show}, \texttt{show-active},
\texttt{show-policies} und \texttt{show-attributes}; \texttt{set} oder
\texttt{set-stream} existieren nicht.
\begin{lstlisting}
sudo ./target/release/tails-pdp-admintool show-active
sudo ./target/release/tails-pdp-admintool show-policies
sudo ./target/release/tails-pdp-admintool show-attributes
./target/release/tails-pdp-admintool --help
\end{lstlisting}
Root ist in den Beispielen für den Map-Zugriff vorgesehen; tatsächlich ausreichende
Leserechte können auch ohne \texttt{sudo} genügen. Optionen wie
\texttt{--policy-dir} und \texttt{--attributes-dir} dienen der Zuordnung der
Quellverzeichnisse bei der Anzeige; sie ändern nicht die Verzeichnisse der Runtime.
Map-Pfade lassen sich über die im CLI definierten Pin-Pfadoptionen auswählen.

\begin{lstlisting}
sudo env RUST_LOG=debug ./target/release/tails-pdp
sudo env TAILS_PDP_EBPF_DEBUG=1 RUST_LOG=info ./target/release/tails-pdp
sudo cat /sys/kernel/debug/tracing/trace_pipe
sudo env TAILS_PDP_TIMING=1 ./target/release/tails-pdp
\end{lstlisting}
\texttt{RUST\_LOG} steuert Userspace-Logs; Standard ist \texttt{info}.
\texttt{TAILS\_PDP\_EBPF\_DEBUG} aktiviert Kernel-Debugausgaben.
Der Trace-Aufruf benötigt eine verfügbare Tracefs-/Debugfs-Ablage.
\texttt{TAILS\_PDP\_TIMING=1} aktiviert opt-in Messmarkierungen auf Basis der
Linux-Monotonzeit. Logs und Marker beeinflussen Performance-Messungen und müssen
bei deren Interpretation berücksichtigt werden.

\section{Build, Tests und Erweiterung}
\subsection{Entwicklungsbefehle}
Buildwerkzeuge und Linux-Voraussetzungen sind im Installationshandbuch beschrieben.
\path{nix-config/rust-toolchain} legt aktuell Rust \texttt{1.93.0} für Userspace fest;
\path{nix-config/ebpf-toolchain} legt \texttt{nightly-2026-01-15} für den eBPF-Build fest.
\path{tails-pdp/build.rs} liest die eBPF-Version und übergibt sie an Aya.
Die Nix-Shell wählt die Userspace-Version und unterbindet implizite Rustup-Installationen.
Der aktuelle Aya-eBPF-Build setzt weiterhin Nightly und dessen \texttt{rust-src} voraus.

Die Quellenersetzung unter \path{.cargo/config.toml} und der Offline-Modus sind aktiv.
\path{vendor/} enthält die Projektabhängigkeiten und die unter
\path{nix-config/ebpf-sysroot/} beschriebenen Registry-Pakete des
Systembibliotheksbuilds. \path{scripts/check-vendor.py} prüft deren Vollständigkeit
und Prüfsummen; der Buildhelfer vergleicht auch das Lockfile der installierten Rust-Quellen.
Compiler, LLVM, Linker und Nix-Store-Abhängigkeiten sind gesonderte Voraussetzungen.

Innerhalb der vorbereiteten Nix-Shell:
\begin{lstlisting}
bash scripts/demo/build.sh
cargo fmt --all -- --check
cargo test --locked -p tails-pdp-common -p tails-pdp-policy-loader --config 'target."cfg(all())".runner="env"'
cargo doc --locked --no-deps -p tails-pdp-common -p tails-pdp-policy-loader
\end{lstlisting}
Der Buildhelfer verwendet \texttt{cargo build --frozen} und prüft vorhandene
Toolchains, Rust-Quellen und Linker. Er installiert keine fehlenden Werkzeuge.
Die Test-Runner-Überschreibung verhindert Root-Ausführung normaler Rust-Tests;
das Workspace-Cargo-Setup verwendet sonst \texttt{sudo -E}.
Der Hauptbuild baut über \path{tails-pdp/build.rs} auch das eBPF-Objekt.
Die Buildskripte verbleiben in ihren Crates. Generierte Kerneltypen werden nur
gezielt untersucht und nicht pauschal verändert.

\subsection{Geführte Demo und Isolation}
\path{scripts/demo.sh} ist der Einstieg für die Demonstration. Der Launcher richtet
tmux ein und startet den privilegierten Worker. Der fachliche Ablauf liegt in
\path{scripts/demo/runner.sh}, die System-sudo-Auswahl in
\path{scripts/demo/dependencies.sh}, der Offline-Build in
\path{scripts/demo/build.sh} und die Terminaldarstellung in \path{scripts/demo/pane.sh}.
\begin{lstlisting}
./scripts/demo.sh
./scripts/demo.sh --auto
./scripts/demo.sh --build
TAILS_PDP_BIN=/absoluter/pfad/tails-pdp ADM_TOOL_BIN=/absoluter/pfad/tails-pdp-admintool ./scripts/demo.sh
\end{lstlisting}
Die Optionen steuern Eingabepausen beziehungsweise einen erzwungenen Neubuild.
Ohne Binaries wird automatisch ein Build angefordert. Fehlende Werkzeuge führen
entweder zur lokalen Nix-Demo-Shell oder zu einer Fehlermeldung; Nix-Downloads
und Nix-Builds sind für diese Bereitstellung deaktiviert. Die vollständige
Offline-Werkzeugumgebung ist derzeit nicht im Repository enthalten.

Ein privater Mount-Namespace mit eigenem bpffs erhält die Host-Pins. Die Demo
sichert vorhandene Konfigurationsverzeichnisse und \path{demo/files/} und stellt
sie beim regulären Aufräumen wieder her. Logs bleiben unter \path{demo/runs/}.
LSM-Hooks wirken weiterhin systemweit: keine zweite Runtime und keine parallele
Evaluation starten. \texttt{DEMO\_TIMEOUT\_SECONDS} begrenzt technische Wartephasen
standardmäßig auf 30 Sekunden. Weitere Optionen dokumentiert \path{scripts/demo/README.md}.

\subsection{API-Dokumentation mit Rustdoc}
Die Crates und ihre öffentlichen Schnittstellen sind unmittelbar im Quelltext mit
\texttt{//!}- bzw. \texttt{///}-Kommentaren dokumentiert. Die Referenz erläutert
Parametersemantik, Rückgaben, Fehlerfälle, Wildcards, Datenlayouts und relevante
Plattformgrenzen. Wesentliche interne Algorithmen sind zusätzlich an den betreffenden
Routinen dokumentiert. Signaturen und Typverweise erzeugt Rustdoc aus dem Code.

Auf dem vorbereiteten Linux-Zielsystem erzeugt ein gemeinsamer Einstieg die
Referenz aller acht Bibliotheksziele:
\begin{lstlisting}
./scripts/docs.sh
./scripts/docs.sh --document-private-items
cargo test --locked --doc -p tails-pdp-common --config 'target."cfg(all())".runner="env"'
\end{lstlisting}
Das Skript prüft fehlende öffentliche Dokumentation und ungültige Rustdoc-Verweise
als Fehler. Die zweite Variante enthält zusätzlich interne Symbole. Die
HTML-Ausgabe liegt unter \path{target/doc/}; bei einem expliziten Cross-Target unter
\path{target/<target>/doc/}. Die gesamte HTML-Ablage einschließlich Ressourcen
kann separat in das Abgabepaket kopiert werden. Sie wird nicht als Teil des
Buildverzeichnisses versioniert. Der Beispieltest benötigt keine BPF-Pins und
keine Root-Rechte; privilegierte Runtime-Tests bleiben ein eigener Nachweis.

\subsection{Testzuordnung}
\begin{lstlisting}
bash scripts/tests/test-publish-config.sh
python3 -B scripts/tests/test_demo.py
PYTHONPATH=scripts/tests/evaluation python3 -B -m unittest discover -s scripts/tests/evaluation -p test_perf03.py
./scripts/test.sh
sudo ./scripts/tests/test-e2e.sh
sudo ./scripts/tests/test-evaluation.sh E2E-18 E2E-19
\end{lstlisting}
Die ersten drei Befehle sind lokale Hilfsprüfungen. Die vollständige Prüfkette,
Rustkomponenten mit Linux-APIs sowie die privilegierten Tests benötigen das
dedizierte Linux-Zielsystem. E2E-01 bis E2E-19 prüfen die reale Runtime.
COMP-, PERF-, LOAD-, STAB-, CHAR- und RACE-Szenarien untersuchen weitere
Komponenten-, Mess- und Grenzfälle. Die vollständige Prüfkette enthält außerdem
\path{scripts/tests/test_demo_live.py}: Der Test führt die automatische Demo mit
realem tmux und Kernel-Enforcement unter einem Pseudoterminal aus und archiviert
Report, Terminalmitschnitt und Demo-Fixture unter \path{test-results/}.
Der privilegierte Live-Test ist kein Ersatz für lokale Demo-Regressionsprüfungen.
Die Suite vor dem Start von fremden
Runtime-Prozessen und vorhandenen Test-Pins freihalten; Suiten nicht parallel starten.

\subsection{Erweiterungsschritte}
\begin{enumerate}
  \item Fachliche Semantik und Fehlerverhalten festlegen. Für zusätzliche Syntax
    Parser und Validierung in \path{tails-pdp-policy-loader/src/policy_source.rs}
    oder \path{tails-pdp-attribute-loader/src/stream_attributes.rs} anpassen.
  \item Gemeinsame Repräsentation und reine Evaluierung in
    \path{tails-pdp-common/src/lib.rs} erweitern. Bestehende Map-Layouts und
    Enumwerte beachten; bei inkompatiblen Änderungen Migration und Pin-Bereinigung
    planen. Kapazitäten nicht nur im Parser, sondern auch in Maps und Prüfungen ändern.
  \item Kernel- und Userspace-Auswertung auf dieselbe Semantik bringen.
    Fehlende Attribute, Typkonflikte, Wildcards und Zeitgrenzen ausdrücklich prüfen.
    Für einen neuen Hook zusätzlich Requestgewinnung, Programme, Maps,
    Lade-/Attach-Listen, Administration und Nachbewertung entwerfen.
  \item Positive und negative Parser-/Komponententests sowie passende
    privilegierte Szenarien ergänzen. Dokumentation und Beispiele aktualisieren;
    erst danach Build, Clippy, Kernel-Verifier und reale Durchsetzung prüfen.
\end{enumerate}
Die Wiederaufnahme einer Socket-Bind-Aktion ist eine neue Funktion und kein
Umschalten eines vorhandenen Flags. Die vorherigen Implementierungsreste sind entfernt.

\section{Grenzen und Nachweiszuordnung}
Der Prototyp ist keine vollständige SAPL-Implementierung. Der Remote-FD-Entzug
ist nur für x86\_64-Linux implementiert und benötigt geeignete ptrace-Rechte.
Bereits erzeugte Speicherabbildungen und konkurrierende FD-/Prozessoperationen
sind nicht durch einen atomaren globalen Entzug abgedeckt. Ein bestandener
Race-Test beweist keine Race-Freiheit. Dateiidentität und Pfadname dürfen nicht
gleichgesetzt werden. Mehrdatei-Veröffentlichungen und Policy-/Attributaktivierung
bilden keine gemeinsame Transaktion. Eine einzelne Generationsumschaltung belegt
keine unveränderte Lesesicht bei wiederholter Bankverwendung durch parallele Leser.
Der synchrone Userspace-Scan und blockierendes Warten beim \texttt{ptrace}-Entzug
können andere Zweige desselben Runtime-Tasks verzögern. Für das Warten besteht
keine ausdrückliche Zeitbegrenzung; eine feste maximale Aktualisierungslatenz wird
nicht zugesichert. Fehlgeschlagene Entzüge werden nicht ohne neuen Auslöser wiederholt.

Das Archiv \path{eval-tails-pdp/} enthält den Lauf vom 10. Oktober 2026 zum
protokollierten Commit \texttt{a9fd6c0} mit dokumentiertem Arbeitsbaum.
Seine Dateien sind über \path{eval-tails-pdp/SHA256SUMS} prüfbar.
Bereinigungen nach diesem Lauf und dieses Handbuch sind kein neuer vollständiger
Systemtestnachweis. Für die Abgabe den endgültigen Quellstand, Binaries und
Nachweise eindeutig zuordnen.

\section{Ausgabe des Handbuchs}
Diese LaTeX-Datei benötigt keine eingebundenen Thesis-Dateien, Grafiken oder
Literaturdaten. Der gemeinsame Build benötigt die Quellen unter \path{../thesis/}
sowie LaTeX, \texttt{latexmk} und Biber. Aus der Repository-Wurzel:
\begin{lstlisting}
bash scripts/build-thesis.sh
\end{lstlisting}
Die lokale Laufkonfiguration \texttt{thesis} verwendet denselben Helfer. Er baut
Thesis und beide Handbücher und kopiert alle drei PDFs erst nach erfolgreicher
Kompilierung nach \path{thesis/final/}. Dieses Handbuch liegt dort als
\path{Christian_Koch_Programmierhandbuch_tails-pdp.pdf}. Die bearbeitbare Quelle
verbleibt für die Abgabe im vorgesehenen Thesis-/source-Paket.
Eine eigenständige Kompilierung dieses Handbuchs bleibt möglich.
\end{document}
